% ============================================================
% Section 133: 质量作用定律与金属-半导体接触
% ============================================================
\section{半导体物理：质量作用定律与金属-半导体接触}
\label{sec:mass-action-schottky}

\begin{modelbox}[题目：硅的质量作用定律与接触特性]
在晶态硅中，能隙是 $1.14\,\text{eV}$，空穴有效质量 $m_h\approx0.3m$，电子有效质量 $m_e=0.2m$。

(1) 作出合适的近似（并证实），求出质量作用定律 $np=f(T)$ 中的函数 $f(T)$ 的表达式。

(2) 要掺入多大浓度的施主杂质砷（As），才能使室温下非本征电导率比本征电导率高 $10^4$ 倍以上？（忽略受主杂质，静态介电常数 $\varepsilon=11.8$）

(3) 证明：依金属功函数 $\Phi_m$ 和半导体功函数 $\Phi_s$ 的相对大小，金属-半导体结合可以是整流结或欧姆结。
\end{modelbox}

\begin{tipbox}[考点快速分析]
\begin{itemize}
    \item \textbf{质量作用定律}：$np=N_cN_ve^{-E_g/k_BT}$，非简并条件下的普适结果
    \item \textbf{有效态密度}：$N_{c,v}=2(2\pi m_{e,h}k_BT/h^2)^{3/2}\propto T^{3/2}$
    \item \textbf{电导率对比}：$\sigma_i=en_i(\mu_e+\mu_h)$，$\sigma\approx eN_D\mu_e$
    \item \textbf{肖特基结/欧姆结}：由功函数差决定能带弯曲方向
\end{itemize}
\end{tipbox}

\begin{derivationbox}[标准解题过程]
\textbf{(1) 质量作用定律 $np=f(T)$}

在非简并条件下，$n=N_ce^{-(E_c-E_F)/k_BT}$，$p=N_ve^{-(E_F-E_v)/k_BT}$。相乘得
\[
np=N_cN_ve^{-E_g/k_BT}.
\]
代入有效态密度表达式：
\begin{equation}
\boxed{f(T)=4\left(\frac{2\pi k_BT}{h^2}\right)^3(m_em_h)^{3/2}e^{-E_g/k_BT}.}
\end{equation}

\textit{近似证实}：室温下本征硅 $n_i\sim10^{10}\,\text{cm}^{-3}\ll N_c,N_v\sim10^{19}\,\text{cm}^{-3}$，费米能级位于禁带中央附近，$E_c-E_F\approx E_g/2\approx0.57\,\text{eV}\gg k_BT$，非简并近似成立。

\textbf{(2) 所需施主掺杂浓度}

本征电导率 $\sigma_i=en_i(\mu_e+\mu_h)$。掺入施主后 $\sigma\approx eN_D\mu_e$（$N_D\gg n_i$，忽略空穴）。要求 $\sigma/\sigma_i\ge10^4$：
\[
\frac{N_D\mu_e}{n_i(\mu_e+\mu_h)}\approx\frac{N_D}{2n_i}\ge10^4\implies N_D\ge2\times10^4n_i.
\]

计算室温 $n_i$：
\[
n_i=2\left(\frac{2\pi k_BT}{h^2}\right)^{3/2}(m_em_h)^{3/4}e^{-E_g/2k_BT}\approx9.15\times10^8\,\text{cm}^{-3}.
\]
因此
\begin{equation}
\boxed{N_D\ge2\times10^4\times9.15\times10^8\approx1.83\times10^{13}\,\text{cm}^{-3}.}
\end{equation}
这是一个极低的掺杂浓度，实际硅器件掺杂通常远高于此值。

\textbf{(3) 金属-半导体接触：整流结与欧姆结}

以 n 型半导体为例。

\textit{情况一：$\Phi_m<\Phi_s$（欧姆接触）}

接触后电子从金属流向半导体，在半导体表面形成负电荷积累，能带向下弯曲，形成\textbf{反阻挡层}（高电导区）。该区域对电流无整流作用，无论外加电压极性如何电流均可畅通，呈现低阻欧姆特性。

\textit{情况二：$\Phi_m>\Phi_s$（整流接触，肖特基结）}

接触后电子从半导体流向金属，在半导体表面留下正电荷（电离施主），形成耗尽层，能带向上弯曲，形成\textbf{肖特基势垒}。正向偏压（金属接正）时势垒降低，电流大；反向偏压时势垒升高，仅极小漏电流。呈现典型的整流特性。

（对于 p 型半导体，结论的功函数大小关系相反。）
\end{derivationbox}

\begin{formulabox}[答案汇总]
\begin{center}
\begin{tabular}{ll}
\toprule
\textbf{问题} & \textbf{答案要点} \\
\midrule
(1) 质量作用定律 & $f(T)=4\bigl(\dfrac{2\pi k_BT}{h^2}\bigr)^3(m_em_h)^{3/2}e^{-E_g/k_BT}$ \\
(2) 所需 $N_D$ & $\ge1.83\times10^{13}\,\text{cm}^{-3}$（极低浓度） \\
(3) $\Phi_m<\Phi_s$ & 欧姆接触（反阻挡层） \\
(3) $\Phi_m>\Phi_s$ & 整流接触（肖特基势垒） \\
\bottomrule
\end{tabular}
\end{center}
\end{formulabox}

\begin{insightbox}[物理图像]
\begin{itemize}
    \item 质量作用定律是非简并半导体的基石，将载流子浓度乘积与材料本征参数和温度联系起来。
    \item 硅的室温本征载流子浓度约 $10^{10}\,\text{cm}^{-3}$，极微量的掺杂即可显著改变电导率，这是半导体器件工作的基础。
    \item 金属-半导体接触的整流或欧姆特性由功函数差决定。肖特基结是高频和功率器件的重要组成部分；欧姆接触则是器件电极的必要条件。
\end{itemize}
\end{insightbox}
